\pdfbookmark[1]{Abstract}{abstract}
\begin{center}
\Large\spacedlowsmallcaps{Abstract}
\end{center}
\bigskip

\begin{itshape}
Nowadays the huge amount of data generated by digital financial systems, together with the evolution of fraud schemes, make fraud detection a challenging and critical problem.
Since illicit actors employ sophisticated strategies such as fraud rings and cyclic transfers to obscure their trails, traditional heuristic and rule-based detection systems frequently fail to provide an effective solution.

The primary goal of this thesis is to develop a detailed Python implementation of a Quantum Topological Graph Neural Network (QTGNN) \cite{mainarticle}. This model combines graph neural networks, quantum state embedding and topological data analysis to capture higher-order, long-range dependencies and robust topological invariants that are typical in fraud schemes.
Leveraging these features we will also be able to enhance the interpretability of the model. This is extremely important since, in regard to using deep learning models in the real world and particularly in legal settings, the model must be able to explain its decisions \cite{goodman2017european}.

The model will be evaluated on four different transaction datasets: the Elliptic Bitcoin Dataset, Ethereum Phishing Transaction Networks, IBM AMLSim, and PaySim. To balance computational feasibility with practical applicability, the model will be trained on a quantum simulator. Only the final testing phase will be executed on an actual NISQ quantum computer.
\end{itshape}
